\chapter{Introduction}\label{ch:intro}

\section{Topic and motivation}

Generative diffusion models have become the dominant paradigm for
synthesising high-dimensional data: they power state-of-the-art systems
for images, audio, video, and molecular structures \citep{sohl2015deep,
ho2020denoising, song2021scorebased, karras2022elucidating}. The same
construction underlies the image and video generators deployed at scale by the
large technology companies, from latent diffusion \citep{rombach2022latent}
onwards, and reaches beyond media: structure prediction in AlphaFold~3 places
atomic coordinates with a diffusion module \citep{abramson2024alphafold3}. Their
central object is the \emph{score function}, the gradient of the
log-density of progressively noised data,
\begin{equation}
  \score(x, t) \;=\; \nabla_x \log p_t(x),
  \label{eq:intro-score}
\end{equation}
\begin{figure}[t]\centering
\begin{tikzpicture}[
  latent/.style={circle, draw, minimum size=6.4mm, inner sep=0pt, font=\small},
  obs/.style={circle, draw, fill=black!12, minimum size=6.4mm, inner sep=0pt, font=\small},
  fac/.style={rectangle, draw, fill=black, minimum size=1.8mm, inner sep=0pt},
  box/.style={rectangle, draw, rounded corners=1.6pt, align=center,
              inner sep=3.6pt, font=\small, text width=45mm},
  lbl/.style={font=\small},
  >={Stealth}, scale=0.95, every node/.style={transform shape}]

% ---- clean chain -----------------------------------------------------
\foreach \i in {0,1,2,3}{ \node[latent] (a\i) at (1.5*\i, 0) {$a_{\i}$}; }
\foreach \i/\j in {0/1,1/2,2/3}{
  \node[fac] (k\i) at (1.5*\i+0.75, 0) {};
  \draw (a\i) -- (k\i) -- (a\j);
}
\node[lbl, left=2.5mm of a0, align=right] {clean\\chain};
\node[lbl, above=0.6mm of k1] {$\kernel(a_{k}\mid a_{k-1})$};

% ---- channel ---------------------------------------------------------
\foreach \i in {0,1,2,3}{
  \node[obs] (x\i) at (1.5*\i, -1.75) {$x_{\i}$};
  \draw[->, thick] (a\i) -- (x\i);
}
\node[lbl, left=2.5mm of x0, align=right] {noised\\at time $t$};
\node[lbl] at (2.25,-2.75) {$x_k = e^{-t}a_k + \sqrt{\Deltat}\,\varepsilon_k$};

% ---- the two routes and the object they compute ----------------------
\coordinate (j) at (6.1,-1.75);
\node[box] (bp)    at (9.3, 0.15) {exact inference\\[1pt]
  {\footnotesize BP, and EM--BP when $\kernel$ is unknown}};
\node[box] (net)   at (9.3,-3.65) {trained network\\[1pt]
  {\footnotesize denoising score matching}};
\node[box] (score) at (9.3,-1.75) {joint score\\[1pt] $\nabla_x \log p_t(x)$};

\draw[thick] (x3) -- (j);
\draw[->, thick] (j) |- (bp.west);
\draw[->, thick] (j) |- (net.west);
\draw[->, thick] (j) -- (score.west);
\draw[->, thick] (bp.south)  -- (score.north);
\draw[->, thick] (net.north) -- (score.south);
\node[lbl, right=2mm of score, align=left, font=\footnotesize]
  {compared at\\equal data};
\end{tikzpicture}
\caption[What this thesis computes, and the comparison it makes]{The object of
the thesis and the two ways it is obtained. A Markov chain generates a clean
trajectory; the Ornstein--Uhlenbeck channel corrupts every coordinate
independently, so each observation attaches to the latent chain on its own and
the posterior stays a chain. That is what makes the joint score exactly
computable by message passing, with the transition kernel known
(Chapters~\ref{ch:gaussian} and~\ref{ch:laplace}), or estimated from noised
sequences alone when it is not (Chapter~\ref{ch:em-method}). The same score is
what a denoising-score-matching network is trained to approximate without being
given the chain structure. Chapter~\ref{ch:em-results} measures what supplying
that structure is worth, at equal training data, and where it stops paying.}
\label{fig:thesis-pipeline}
\end{figure}

which is the one object a model must learn in order to reverse a noising
process: given the forward process, its schedule and its terminal law, the score
determines the reverse drift \citep{anderson1982reverse}. In
practice the score is approximated by a neural network trained by
denoising score matching \citep{hyvarinen2005estimation,
vincent2011connection}; the network is a black box, and the structure of
the function it has to learn remains poorly understood.

A complementary research programme, initiated in the statistical-physics
community, replaces the black box with \emph{exactly solvable models}:
data distributions simple enough that the score can be computed in
closed form, yet rich enough to exhibit the phenomena observed in real
systems, such as speciation, memorisation, and dynamical phase
transitions along the generation trajectory \citep{biroli2023generative,
biroli2024dynamical, raya2023spontaneous, ghio2024sampling,
achilli2026speciation}. Exact scores act as a microscope: every regime,
transition time, and failure mode can be traced back to explicit
formulas.

This thesis works inside that programme, in a direction that is natural
for data with \emph{temporal or sequential structure}: the data are not
a cloud of independent points but a \emph{trajectory}
$a_0, a_1, \dots, a_{\nsites-1}$ of a Markov process, the prototype of a
video, a time series, or any dynamic object. Each frame of the
trajectory is corrupted independently by an Ornstein--Uhlenbeck (OU)
channel, and the object of study is the \emph{joint} score of the whole
noisy sequence. Two classical bodies of theory then meet the modern one.
First, the joint score of a Gaussian Markov chain is governed by the
algebra of structured matrices: Toeplitz covariances, tridiagonal
precisions, and their common eigenbasis. Second, by an identity that
holds for any data distribution, computing a score is solving a Bayesian
inference problem, and for chain-structured data that problem lives on a
tree-shaped factor graph, the domain of \emph{belief propagation} (BP)
and of the Kalman filter, together with the Gaussian message closures
of statistical physics (TAP, AMP)
\citep{pearl1988probabilistic, mezard2009information, kalman1960new,
donoho2009message}. The thesis develops both routes in full, proves
where they coincide, and measures what is lost when they are forced
apart.

\section{Positioning and gap}

Chapter~\ref{ch:statmech} reviews the literature this work builds on; the gap
it leaves is stated here.

Three strands meet. The statistical physics of generative diffusion has
produced closed-form analyses of the backward dynamics and their transitions
\citep{biroli2023generative, biroli2024dynamical, achilli2026speciation,
sclocchi2025phase, bonnaire2025memorize}. Exact posterior inference on
Gauss--Markov chains is classical, through the Kalman filter and smoother and
their formulation as sum--product message passing
\citep{kalman1960new, rauch1965maximum, kschischang2001factor,
mezard2009information}. And message passing with continuous variables closes
its messages inside a parametric family, which is the step connecting belief
propagation to AMP and to the TAP equations
\citep{thouless1977solution, donoho2009message, zdeborova2016statistical}.

Each strand stops short of the same place. The physics analyses treat data as
unstructured collections of samples: the only time axis they carry is the
diffusion time, and the \emph{internal} time of a sequence, a frame correlated
with its neighbours by a dynamical rule, does not appear. The inference
literature has the machinery for that internal time but does not ask what
generative modelling forces it to ask, which is how the exact posterior, and
so the exact score, deforms as the noise level is swept from zero to
saturation. And the standard justification for closing messages in a Gaussian
family is a central-limit argument that requires each node to receive many
messages; on a chain each node receives exactly two, so the argument fails by
construction and we are not aware of a quantitative assessment of what the
closure costs on chain-structured diffusion posteriors.

The gap this thesis addresses is the intersection: the joint score of a noised
Markov sequence, computed exactly where that is possible, approximated in a
controlled and measured way where it is not, and learned from noised data when
the dynamics are unknown.

One methodological choice follows from the same reasoning and is worth stating
early. Where no closed form exists, the reference used throughout is
deterministic quadrature and not Monte Carlo. A particle smoother carries
sampling error, which would add a second error term to every comparison and
force the cost of an approximation and the cost of the reference to be
disentangled before either could be read. Grid-based inference carries a
discretisation error instead, which can be bounded separately and reported,
and that is what makes it usable as a reference and not merely as another
competitor.

\section{Research questions}\label{sec:rq}

The thesis addresses six questions. The first four concern scores under
\emph{known} dynamics and the last two concern learning dynamics that are
unknown.

\begin{itemize}
  \item[\textbf{RQ1.}] For a stationary Gaussian AR(1) chain observed
  through coordinatewise Ornstein--Uhlenbeck noise, what is the joint
  score, and how do its spatial couplings vary with the diffusion time?

  \item[\textbf{RQ2.}] How can the same Gaussian-chain score be computed
  by belief propagation, and how does the resulting recursion relate to
  Kalman filtering and smoothing?

  \item[\textbf{RQ3.}] For selected non-Gaussian innovation models, how
  accurate is a Gaussian-message approximation relative to a validated
  numerical belief-propagation reference, over the parameter ranges
  studied?

  \item[\textbf{RQ4.}] In the Gaussian model, how rapidly does the
  influence of remote observations on the score decay, and what limited
  guidance does this provide for local approximations of the score?

  \item[\textbf{RQ5.}] Under what assumptions can an \emph{unknown} local
  transition law be recovered from noised sequences alone, and which
  numerical, optimisation and statistical checks are required before such
  a recovery claim is reliable?

  \item[\textbf{RQ6.}] On a family where the Markov assumption holds
  exactly, what is supplying that assumption worth, measured against
  denoising-score-matching networks trained on identical data and given
  none of it?
\end{itemize}

These questions are deliberately confined to the models named in them.
No claim is made about general Markov processes, about universal
properties of learned diffusion scores, or about practical neural
architectures. In particular RQ6 is a question about the value of a
correct structural assumption on one family, not a comparison of learning
algorithms in general: the two arms are asymmetric in supplied
information by construction, which is what makes the measurement
meaningful and also what bounds it.

Chapter~\ref{ch:ring} answers no numbered question. It comes first because it
is where the work began, and because what it settles motivates the six: on a
rotating-ring model, whose state space and estimand both differ from the
chain's, it shows what joint observations can identify that no single-frame
marginal can.

\section{Contributions}

\noindent Four contributions, in the order the body develops them. Each is
stated here as a claim and a kind of claim; the derivations are in the chapters
named and the measured values are reported there and in
Chapter~\ref{ch:conclusions}.

\begin{enumerate}
  \item \textbf{The joint score of a noised Gaussian Markov chain, in closed
  form.} For the stationary AR(1) chain under coordinatewise OU corruption we
  derive the score, diagonalise it in a single eigenbasis valid at every
  diffusion time, and give the law by which its couplings fill in at small $t$
  and decay back to isotropy at large $t$, together with exact bulk formulas
  for the posterior variance and the correlation decay rate $q(\alpha,t)$
  (Chapter~\ref{ch:gaussian}). Computing the same score by message passing, we
  show that belief propagation on the chain is the Kalman filter followed by
  Rauch--Tung--Striebel smoothing, update by update, and prove that the
  influence of evidence at distance $d$ falls as $q(\alpha,t)^d$. That the
  radius-$r$ windowed estimator's error falls at the same rate is established
  separately, and by measurement, because windowing changes the boundary
  conditions of the recursion in a way the influence argument does not cover
  (Chapter~\ref{ch:gaussian-bp}).

  \item \textbf{An exact reading of what Gaussian message closure computes.}
  For Laplace innovations the functional recursion stays exact on the chain but
  its messages no longer close in any finite-dimensional family. We build a
  numerically validated quadrature reference, and prove that moment-projected
  messages return exactly the LMMSE estimator of the covariance-matched
  Gaussian model (Theorem~\ref{thm:l-lmmse}). The measured departure is
  therefore attributable: it is what discarding higher moments costs, for every
  innovation law with the same first two moments, and not what one
  implementation costs (Chapter~\ref{ch:laplace}).

  \item \textbf{Recovery of an unknown transition law from noised sequences.}
  Dropping the assumption that the kernel is known, we show that Fisher's
  identity makes one forward--backward pass yield the exact gradient of the
  marginal log-likelihood, so no derivative is taken through the recursion and
  the E-step stays structurally exact for any innovation law
  (Chapter~\ref{ch:em-method}). Measuring what convergence costs then produces
  a result rather than a diagnostic: the correlation and the innovation shape
  settle at very different rates, so a stopping rule read off the first
  certifies a kernel unconverged in the second (Chapter~\ref{ch:em-kernel}).

  \item \textbf{The measured value of supplied Markov structure, and its
  limit.} At equal training data the structured estimator is substantially more
  accurate than networks trained by denoising score matching, and remains so
  against a baseline given locality and weight sharing by an architecture
  screen. The contribution is as much the boundary as the margin: controlled
  departures from the Markov assumption locate where the advantage ends and the
  networks overtake it (Chapter~\ref{ch:em-results}).
\end{enumerate}

\section{Structure of the thesis}

Chapter~\ref{ch:statmech} is the context: statistical mechanics and the
variational principle behind the Boltzmann distribution, energy-based models as
the same object with a learned energy, generative diffusion and the reverse-time
SDE, and then the distinction the rest of the thesis turns on, between the
exact, the empirical and the fitted score. It closes with the reading of U-Nets
as belief propagation that makes exact inference a serious alternative to a
trained network rather than a toy one. Chapter~\ref{ch:model} fixes the model
and the estimand: the Markov sequence, the Ornstein--Uhlenbeck channel, the
argument that the joint score and not a per-frame score is the object worth
studying, and the factor-graph language the computations are written in.

Chapter~\ref{ch:ring} begins the research where the research began, with a
planar trajectory near a circle. It shows that the rotation rate carries zero
Fisher information in every one-frame marginal and positive information in the
trajectory. That is the argument for the joint score in a single exact
statement, and it is what sends the following chapters to a model where the
joint score can be computed in full.

Chapters~\ref{ch:gaussian} and~\ref{ch:gaussian-bp} do that for the stationary
Gaussian chain, first with matrices and then with messages: the score in closed
form, the lifecycle of its couplings, the identification of belief propagation
with the Kalman filter and the Rauch--Tung--Striebel smoother, and the
geometric decay that says how local an accurate approximation can be
(RQ1, RQ2, RQ4). Chapter~\ref{ch:laplace} leaves Gaussianity in the most
controlled way available, through Laplace innovations, builds a validated
numerical reference, and measures exactly what a Gaussian message family
discards (RQ3).

Chapters~\ref{ch:em-method} and~\ref{ch:em-kernel} drop the assumption that the
transition law is known and recover it from noised sequences alone, first two
parameters and then a whole kernel, with the convergence schedule of the two
coordinates as a result in its own right (RQ5). Chapter~\ref{ch:em-results}
measures what supplying the Markov structure is worth against networks trained
on identical data, and locates where it stops being worth anything (RQ6).
Chapter~\ref{ch:conclusions} answers the six questions and states what is left
open.

Each concept is developed once, in the chapter where it is used. The
appendices carry the Gaussian identities, the ring derivations, the
reproducibility protocol and the aggregation robustness of the reported ratios.
Extended background, alternative derivations, the AMP/TAP analysis and the
development record are in \compendium{}.
